% !TEX root = ../../tree_notes.tex
% =====================================================================
%  3.4  Rebuilding and Serialising                       sec:c03-rebuild
%  Exports: def:c03-reconstruct  def:c03-serialize  thm:c03-pair-uniqueness
%           cor:c03-serialize-bijection  lst:c03-build  ex:c03-build-run
%           prob:c03-build-in-post  prob:c03-dup-subtrees
% =====================================================================
\section{Rebuilding and Serialising}\label{sec:c03-rebuild}

\begin{previouslybox}
\begin{enumerate}[nosep, leftmargin=1.6em]
  \item Which order starts with the root, and which puts the root at a fixed split
        position? (\thmref{c03-root-position})
  \item On $R$, state the preorder and the inorder. (\secref{c03-dfs}, \figref{c03-running})
  \item Why does inorder of a \emph{general} binary tree not come out sorted, and where
        do we exploit that caveat here? (Trap 3, \secref{c03-dfs})
\end{enumerate}
\end{previouslybox}

\begin{topicintro}
Everything so far mapped \emph{tree $\to$ sequence}. This section runs the arrow
backwards, and its close cousin, \emph{tree $\leftrightarrow$ string}. Given two traversal
orders we rebuild the unique tree (\term{reconstruction}); given a tree we flatten it to a
token string and parse it back (\term{serialisation}). Both hinge on one asymmetry from
\thmref{c03-root-position}: the root is free-standing in preorder/postorder but is a
\emph{splitter} in inorder -- and it is exactly inorder that removes ambiguity. These are
the deepest tools of the chapter, and they power the hardest interview tree questions
(construction, persistence, subtree hashing, and this chapter's Summit, LC 297).
\textbf{Prerequisites.} \secref{c03-dfs} (all three orders and \thmref{c03-root-position});
maps for index lookup.

\textbf{After this section you can}
\begin{itemize}[nosep]
  \item reconstruct the unique binary tree from inorder $+$ (pre- or post-order);       \hfill\lv{4}
  \item explain precisely why preorder $+$ postorder is ambiguous;                      \hfill\lv{4}
  \item serialise a tree to a string and parse it back, and use a canonical serialisation
        to detect duplicate subtrees;                                                   \hfill\lv{5}
\end{itemize}
\end{topicintro}

\begin{recognitionbox}[Pattern card -- rebuild from orders / serialise]
\reckey{``construct the tree given its X and Y traversals'', ``serialise / deserialise a
binary tree'', ``encode and decode'', ``find all duplicate subtrees'' -- the tree must be
\emph{recovered} from, or \emph{round-tripped} through, a linear form.}
\recuse{
  \item two order-sequences are handed to you and the \emph{structure} must be rebuilt;
  \item the tree must be stored / transmitted as a string and later parsed (\chref{c06}
        reuses this as flat storage);
  \item subtree \emph{identity} is decided by comparing a canonical string (hashing).
}
\recnotuse{
  \item you only need to \emph{emit} an order -- \secref{c03-dfs}, no reconstruction;
  \item the answer is per-row -- \secref{c03-bfs};
  \item the tree is a BST and values must stay ordered -- construction there leans on the
        BST invariant, \chref{c05} (LC 449 is the BST specialisation).
}
One-line test: \emph{is the tree itself the output (or must a string describe it exactly)?
If yes, reconstruction/serialisation.}
\end{recognitionbox}

\subsection{Intuition}

Take $R$'s two orders. Preorder \emph{begins} at the root ($1$); inorder \emph{places} the
root between its two subtrees. Combine them and the whole tree falls out in one stroke:
read the root from preorder, locate it in inorder, and inorder's left part \emph{is} the
left subtree while its right part \emph{is} the right subtree -- recurse.

\begin{figure}[H]
\centering
\begin{tikzpicture}[diagram, font=\small]
  \node[annot, anchor=west] at (-6.6,0.6) {$\mathit{pre}=1,2,4,5,3,6$};
  \node[annot, anchor=west] at (-6.6,-0.1) {$\mathit{in}=4,2,5,\mathbf{1},3,6$};
  \node[annot, text=accentorange, anchor=west] at (-6.6,-0.9)
    {root $=1$ splits inorder into};
  \node[annot, text=primaryblue, anchor=east] at (-2.6,-1.6) {\{4,2,5\} left};
  \node[annot, text=deepblue, anchor=west] at (-1.9,-1.6) {right \{3,6\}};
  \draw[<->, primaryblue] (-6.3,-0.35) -- (-4.9,-1.35);
  \draw[<->, deepblue]  (-0.6,-0.35) -- (-1.3,-1.35);
\end{tikzpicture}
\caption{Reconstruction pivot: preorder names the root ($1$); inorder then splits into the
left subtree's inorder ($4,2,5$, size $3$) and right subtree's inorder ($3,6$). The left
size tells us the next $3$ preorder tokens ($2,4,5$) belong to the left subtree and the
rest to the right -- recurse on both.}\label{fig:c03-rebuild}
\end{figure}

Why is \emph{inorder} indispensable? Without a splitter, a lone child is unattributable.
Preorder and postorder \emph{both} agree on $1,2$ / $2,1$ whether $2$ is the left or the
right child of $1$ -- so (pre, post) cannot tell a one-child node's side:

\begin{figure}[H]
\centering
\begin{tikzpicture}[diagram, scale=0.75, transform shape]
  \node[troot] (a) at (0,0) {$1$};
  \node[tleaf] (b) at (-1,-1.3) {$2$};
  \draw[tedge] (a) -- (b);
  \node[annot, anchor=north] at (0,0.55) {left child};
  \node[troot] (c) at (4,0) {$1$};
  \node[tleaf] (d) at (5,-1.3) {$2$};
  \draw[tedge] (c) -- (d);
  \node[annot, anchor=north] at (4,0.55) {right child};
  \node[annot] at (2,-0.65) {both give\\$\mathit{pre}=1,2$\\$\mathit{post}=2,1$};
\end{tikzpicture}
\caption{The (preorder, postorder) ambiguity: two genuinely different trees share both
orders. Only a node with \emph{exactly one} child is ambiguous; if every node has $0$ or
$2$ children the pair does determine the tree. Inorder resolves the general case.}\label{fig:c03-ambig}
\end{figure}

\subsection{Definitions and Notation}

\begin{definition}{Reconstruction from a traversal pair}{c03-reconstruct}
Given the inorder and one of \{preorder, postorder\} of a tree with \emph{distinct}
values, \term{reconstruction} recovers the unique binary tree: the second order supplies
the root, the inorder supplies the left/right split (and split sizes), and each half is
reconstructed recursively.
\end{definition}

\begin{definition}{Serialisation (with null markers)}{c03-serialize}
A \term{serialisation} is the string obtained by a preorder walk that emits each real node's
value \emph{and} emits a distinguished marker (say \code{\#}) for every \nil{} child, with
an explicit separator. \term{Deserialisation} parses that string by the same preorder,
consuming tokens left-to-right and rebuilding children recursively.
\end{definition}

\begin{vsbox}{Two orders vs a marked string}{reconstruct (pre+in)}{serialise (pre+\#)}
{ Uniqueness needs a \emph{splitter} (inorder) plus distinct values; you are given two
arrays and must interleave them. }
{ Uniqueness is bought instead by writing \emph{null markers} so the missing-child
ambiguity of \figref{c03-ambig} is made explicit -- one walk, self-delimiting. }
\end{vsbox}

\subsection{Properties}

\begin{theorem}{Pair Uniqueness}{c03-pair-uniqueness}
For binary trees with distinct node values: (preorder, inorder) determines the tree
uniquely, as does (postorder, inorder). (Preorder, postorder) does \emph{not} in general.
\end{theorem}
\begin{proofbox}
\emph{(pre, in):} by \thmref{c03-root-position} the first preorder token is the root $r$;
in inorder $r$ occurs at index $k$, so $\mathit{in}[0..k{-}1]$ is exactly
$\mathit{in}(L)$ and $\mathit{in}[k{+}1..]$ is $\mathit{in}(R)$, with $\size(L)=k$. Since
values are distinct, $r$'s inorder position is unique, so the split is forced; the first
$k$ preorder tokens (after $r$) are $\mathit{pre}(L)$ and the remainder $\mathit{pre}(R)$.
Recurse; the base case (empty inorder range) yields $\nil$. Every choice is forced, giving
uniqueness; the same argument reads the root from the \emph{last} postorder token for
(post, in). \emph{(pre, post) fails:} \figref{c03-ambig} exhibits two distinct trees with
identical preorder and postorder -- the side of a single child is unrecoverable because
neither order separates $L$ from $R$ without inorder.
\end{proofbox}

\begin{corollary}{Serialisation is a Bijection}{c03-serialize-bijection}
For distinct (or arbitrary) values, the null-marked preorder of \defref{c03-serialize} is
$\mathit{tree} \leftrightarrow \mathit{string}$ one-to-one: distinct trees yield distinct
strings, and every well-formed string parses to exactly one tree.
\end{corollary}
\begin{proofbox}
The marked preorder is precisely the tree read in preorder with \code{\#} recording each
missing child, so it reconstructs by the deterministic recursive parse (read a token: if
\code{\#} return $\nil$, else make a node and parse its left then right). Because the walk
visits the structure in fixed order, two different trees differ at the first node whose
child-presence or value differs, hence differ in the string -- injective; the parse inverts
it, giving the bijection. (Note the contrast with \thmref{c03-pair-uniqueness}: markers
resolve the one-child ambiguity that (pre,post) cannot.)
\end{proofbox}

\subsection{Reconstruction Algorithm}

\begin{lstlisting}[caption={Rebuild from inorder + postorder (indices, with a value$\to$inorder-index map)}, label={lst:c03-build}]
TreeNode* build(const vector<int>& in, const vector<int>& post,
                int il, int ir, int pl, int pr, unordered_map<int,int>& inIdx) {
    if (pl > pr) return nullptr;                   // empty range -> empty subtree
    int rootVal = post[pr];                        // postorder's last is the root
    TreeNode* root = new TreeNode(rootVal);
    int k = inIdx[rootVal];                        // root's inorder position
    int leftSize = k - il;                         // # nodes in the left subtree
    root->left  = build(in, post, il, k - 1, pl, pl + leftSize - 1, inIdx);
    root->right = build(in, post, k + 1, ir, pl + leftSize, pr - 1, inIdx);
    return root;
}
// caller: fill inIdx[x]=index, then build(in, post, 0, n-1, 0, n-1, inIdx)
\end{lstlisting}
\begin{complexitybox}
With the inorder index map, the root's split position is $\bigO{1}$, so the recurrence is
$T(n) = 2T(n/2) + \bigO{1}$ for a balanced split $=\bigO{n}$; building the map is
$\bigO{n}$. Without it (linear \code{find} each level) you get $\bigO{n^2}$ on a skewed
tree -- the map is not a micro-optimisation, it is what keeps construction linear.
\end{complexitybox}

\subsection{Worked Example}

\begin{example}{Reconstructing $R$ from inorder $+$ postorder}{c03-build-run}
With $\mathit{in}=4,2,5,1,3,6$ and $\mathit{post}=4,5,2,6,3,1$: root $=1$ (last of post),
inorder index $3$ $\Rightarrow$ left size $3$. Left half reads $\mathit{in}[0..2]=4,2,5$,
$\mathit{post}[0..2]=4,5,2$ giving root $2$ with left $4$, right $5$; right half
$\mathit{in}[4..5]=3,6$, $\mathit{post}[3..4]=6,3$ giving root $3$ with \emph{left size
$0$} and right child $6$ -- the skewed node $3$ from \figref{c03-running} is rebuilt
correctly because inorder told us $6$ lies to $3$\textquotesingle s \emph{right}.
\end{example}
\begin{tracebox}[Index windows per recursive call]
\begin{itemize}[nosep, leftmargin=1.4em]
  \item build $[0..5]\mid[0..5]$: root $1$, $k{=}3$, leftSize$3$
  \item left $[0..2]\mid[0..2]$: root $2$, $k{=}1$, leftSize$1$
    \begin{itemize}[nosep,leftmargin=1.6em]
      \item $[0..0]\mid[0..0]$ root $4$; $[2..2]\mid[2..2]$ root $5$
    \end{itemize}
  \item right $[4..5]\mid[3..4]$: root $3$, $k{=}4$, leftSize$0$ $\Rightarrow$ left
        \nil, right $[5..5]\mid[3..3]$ root $6$
\end{itemize}
\end{tracebox}

\subsection{Variations and Twists}

\begin{itemize}[nosep]
  \item \textbf{pre + in} instead of post + in: identical, but take the root from
        \emph{front} of preorder. (LC 105 -- the mirror of \probref{c03-build-in-post}.) \hfill\lv{4}
  \item \textbf{pre + post}: only yields \emph{a} valid tree, not the unique one; build
        left greedily. (LC 889, \secref{c03-rebuild} practice.)        \hfill\lv{3}
  \item \textbf{Serialise the whole tree}: null-marked preorder \emph{is} a
        deserialisable code -- the Summit, LC 297 (chapter end).        \hfill\lv{5}
  \item \textbf{Subtree fingerprint}: a canonical serialisation identifies a subtree; count
        them to find duplicates (\probref{c03-dup-subtrees}).          \hfill\lv{5}
\end{itemize}

\subsection{Common Pitfalls}\label{sec:c03-rebuild-pitfalls}

\begin{warningbox}[Trap 1 -- linear search for the root split]
Calling \code{find(in.begin(), in.end(), rootVal)} at every recursion re-scans the inorder
array, giving $\bigO{n^2}$ and TLE on skewed inputs. Precompute an
\code{unordered\_map<int,int>} value $\to$ index once (\coderef{c03-build}).
\end{warningbox}
\begin{lstlisting}[caption={\textcolor{warningred}{WRONG}: $\bigO{n}$ split-search per node}, language=C++, basicstyle=\ttfamily\small]
int k = find(in.begin()+il, in.begin()+ir+1, rootVal) - in.begin();  // O(n) each
\end{lstlisting}
\begin{lstlisting}[caption={\textcolor{successgreen}{RIGHT}: one-time index map, $\bigO{1}$ split}, language=C++, basicstyle=\ttfamily\small]
// precompute: for (i) inIdx[in[i]] = i;
int k = inIdx[rootVal];                                              // O(1)
\end{lstlisting}

\begin{warningbox}[Trap 2 -- serialising without null markers]
\code{"1,2"} as preorder is ambiguous (is $2$ left or right of $1$? -- \figref{c03-ambig}).
Emit a marker for every \nil{} child so the string is self-delimiting and deserialises to
exactly one tree.
\end{warningbox}
\begin{lstlisting}[caption={\textcolor{warningred}{WRONG}: bare preorder -- cannot decode one-child nodes}, language=C++, basicstyle=\ttfamily\small]
void ser(TreeNode* t, string& s){ if(!t) return; s+=to_string(t->val)+",";
    ser(t->left,s); ser(t->right,s); }         // nulls invisible
\end{lstlisting}
\begin{lstlisting}[caption={\textcolor{successgreen}{RIGHT}: mark the nulls}, language=C++, basicstyle=\ttfamily\small]
void ser(TreeNode* t, string& s){ if(!t){ s+="#,"; return; }   // emit marker
    s+=to_string(t->val)+","; ser(t->left,s); ser(t->right,s); }
\end{lstlisting}

\begin{warningbox}[Trap 3 -- no separator between values]
For values $\ge 10$, \code{"123"} could be $12,3$ or $1,23$; even for single digits a
negative sign needs care. Always join tokens with a delimiter (\code{,}) and split on it
when parsing.
\end{warningbox}
\begin{lstlisting}[caption={\textcolor{successgreen}{RIGHT}: delimited tokens survive multi-digit and negative values}, language=C++, basicstyle=\ttfamily\small]
// "-3,#,#," parses unambiguously to a single left node holding -3
\end{lstlisting}


\subsection{Problem Lab}

Both are \lv4--\lv5 constructions: the first rebuilds a tree from a traversal pair (the
\coderef{c03-build} template, fully deployed), the second turns \emph{subtrees} into
strings to detect repeats -- a serialisation used as an identity test, and a warm-up for
this chapter's Summit (LC 297).

\solvedproblem{c03-build-in-post}{Construct Binary Tree from Inorder and Postorder Traversal}{LeetCode 106}
  {https://leetcode.com/problems/construct-binary-tree-from-inorder-and-postorder-traversal/}{LC Medium L4}
  {divide-and-conquer reconstruction, index map}

\begin{problemstatement}[Inorder + Postorder $\to$ Tree]
Given two integer arrays \code{inorder} and \code{postorder} (distinct values) that are
the inorder and postorder traversals of the same binary tree, rebuild and return that
tree. $n \le 3000$.
\end{problemstatement}

\begin{attemptgate}[Attempt it yourself first -- \lv{4}, 25 minutes]
A fair attempt: the postorder's last value is the root; use it to cut the inorder array.
\peek{which array tells you the \emph{size} of the left subtree?}
\nudge{root's inorder index $k$ gives left size $k-il$; that size carves the postorder
range for the left subtree. Recurse on both windows.}
\fullhint{\coderef{c03-build}: read root from \code{post[pr]}, split via an
\code{inIdx} map, recurse. Full build below.}
\end{attemptgate}

\solstep{Understanding the Problem}
We must emit \emph{pointers}, not a number, so the algorithm constructs nodes and wires
them. The information is exactly \thmref{c03-pair-uniqueness}: postorder's tail names the
root; inorder's index of that root fixes the left/right boundary and the left subtree's
\emph{size}, which in turn carves the postorder ranges.

\solstep{Key Observations}
\begin{observation}
Root $= \mathit{postorder}$'s \emph{last} element of the current window (\defref{c03-three-orders}).
\end{observation}
\begin{observation}
If the root sits at inorder index $k$, the left subtree has $k-il$ nodes; those occupy the
next $k-il$ slots of postorder, and the right subtree's postorder window ends one before
the root.
\end{observation}
\begin{observation}
Distinct values make the value $\to$ inorder-index map a function (one index per value),
which is what turns the split into an $\bigO{1}$ lookup.
\end{observation}

\solstep{Think It Through With Me}
\begin{thinkbox}
\deadend{Recurse using \code{find()} to locate the root in inorder.}{Correct output, but
$\bigO{n}$ per node $\Rightarrow$ $\bigO{n^2}$ overall; at $n=3000$ on a skewed tree this
is $\sim 9\cdot10^{6}$ with heavy constants and TLE risk.}
\deadend{Rebuild from postorder alone, guessing splits.}{Postorder without a splitter does
not say where the left subtree ends -- \thmref{c03-pair-uniqueness} says inorder is needed.}
\aha{Precompute \code{inIdx}; then each call is $\bigO{1}$ work plus two recursive calls
on disjoint windows -- total $\bigO{n}$.}
\end{thinkbox}

\solstep{The Approach Ladder}
\begin{center}\renewcommand{\arraystretch}{1.3}
\begin{tabular}{@{}c p{0.235\linewidth} p{0.28\linewidth} l c@{}}
\toprule
\textbf{\#} & \textbf{Approach} & \textbf{Idea} & \textbf{Time} & \textbf{Verdict} \\
\midrule
1 & \lv{3} recursion + \code{find} & linear split search & $O(n^2)$ & \tle \\
2 & \lv{4} recursion + index map & $\bigO{1}$ split & $O(n)$ & \ac \\
\bottomrule
\end{tabular}\end{center}

\solstep{Solution}
\coderef{c03-build} verbatim: build \code{inIdx}, then \code{build(in,post,0,n-1,0,n-1)}.
Each call takes root from the postorder window's tail, reads its inorder index, and
recurses on the two child windows computed from the left size.

\solstep{Proof of Correctness}
\begin{proofbox}
By \thmref{c03-pair-uniqueness} there is exactly one tree with the given (inorder,
postorder). We show \coderef{c03-build} returns it, by induction on the window length.
Base (empty window, $\mathit{pl}>\mathit{pr}$): return $\nil$, correct for the empty
subtree. Step: the last postorder token is the subtree root (def of postorder); its inorder
index $k$ splits inorder into the left/right subtree traversals (def of inorder,
\thmref{c03-root-position}); $\mathit{leftSize}=k-il$ equals the left subtree's node count,
so the postorder window $[\mathit{pl},\mathit{pl}+\mathit{leftSize}-1]$ is exactly the left
subtree's postorder and $[\mathit{pl}+\mathit{leftSize},\mathit{pr}-1]$ the right's. By the
hypothesis both recursive calls rebuild those subtrees correctly; wiring them under the new
root yields the unique tree.
\end{proofbox}

\solstep{Implementation}
\begin{lstlisting}[caption={Construct from Inorder + Postorder (LeetCode 106)}, label={lst:c03-build-in-post}]
#include <bits/stdc++.h>
using namespace std;

struct TreeNode {
    int val;
    TreeNode *left, *right;
    TreeNode(int x = 0) : val(x), left(nullptr), right(nullptr) {}
};

class Solution {
    unordered_map<int,int> inIdx;
    TreeNode* build(const vector<int>& in, const vector<int>& post,
                    int il, int ir, int pl, int pr) {
        if (pl > pr) return nullptr;
        int rootVal = post[pr];
        TreeNode* root = new TreeNode(rootVal);
        int k = inIdx[rootVal];
        int leftSize = k - il;
        root->left  = build(in, post, il, k - 1, pl, pl + leftSize - 1);
        root->right = build(in, post, k + 1, ir, pl + leftSize, pr - 1);
        return root;
    }
public:
    TreeNode* buildTree(vector<int>& in, vector<int>& post) {
        for (int i = 0; i < (int)in.size(); ++i) inIdx[in[i]] = i;
        return build(in, post, 0, (int)in.size() - 1, 0, (int)post.size() - 1);
    }
};
\end{lstlisting}

\solstep{Trace}
\begin{tracebox}[On the running example $R$]
Given $\mathit{in}=[4,2,5,1,3,6]$, $\mathit{post}=[4,5,2,6,3,1]$.
Root $1$ ($k{=}3$, leftSize $3$). Left window $\to$ root $2$ (leftSize $1$) with children
$4,5$; right window $\to$ root $3$ (leftSize $0$, so left $\nil$) with right child $6$.
Rebuilt tree $=R$ (\figref{c03-running}). \ac.
\end{tracebox}

\solstep{Complexity}
\begin{complexitybox}
Time $\bigO{n}$ (each node built with $\bigO{1}$ split work); space $\bigO{n}$ for the map
plus $\bigO{h}$ recursion stack ($\bigO{n}$ skewed). Building the map is $\bigO{n}$.
\end{complexitybox}

\solstep{Verification and Edge Cases}
\begin{center}\renewcommand{\arraystretch}{1.2}
\begin{tabular}{@{}p{0.30\linewidth} p{0.62\linewidth}@{}}
\toprule
\textbf{Case} & \textbf{What the code does} \\
\midrule
Single node & one call, both windows size $1$, leaf returned \\
Skewed (all-left) tree & leftSize $=$ full range each time; still $\bigO{n}$ with the map,
  but $\bigO{n^2}$ without it (Trap 1) \\
Empty tree & initial \code{pl>pr} guard returns $\nil$ \\
\bottomrule
\end{tabular}\end{center}

\solstep{Pitfalls and Variants}
\begin{warningbox}
Off-by-one on the postorder right window (\code{pr-1}, excluding the root) is the most
common construction bug -- draw the four indices on paper for a 3-node tree before coding.
Variant LC 105 (preorder+inorder) is the mirror: root from the \emph{front} of preorder,
windows shifted accordingly. LC 889 (preorder+postorder) must \emph{choose} a side because
the answer is non-unique (\figref{c03-ambig}).
\end{warningbox}

\begin{insightbox}[Key takeaways -- three things to remember]
\begin{enumerate}[nosep, leftmargin=1.6em]
  \item Reconstruction is divide-and-conquer where the \emph{pivot} is the root's inorder
        position -- the splitter is the whole algorithm.
  \item Turning a per-node search into a one-time index map is the difference between
        $\bigO{n^2}$ and $\bigO{n}$; amortise the lookup, not the recursion.
  \item ``Return a \code{TreeNode*} and wire children'' is the constructive form of
        structural recursion from \secref{c02-transform}.
\end{enumerate}
\end{insightbox}

\begin{hardeningbox}
\begin{itemize}[nosep, leftmargin=1.4em]
  \item \textbf{Deep recursion}: $n=3000$ skewed is $\bigO{n}$ stack depth -- fine for LC,
        but a contest may push $n$ higher; an iterative build with an explicit stack of
        windows mirrors \secref{c03-iterative}'s discipline.
  \item \textbf{Duplicate values}: \thmref{c03-pair-uniqueness} needs distinct values; if
        values repeat the map is no longer a function and the tree is not unique -- the
        problem promises distinctness, and you should say so when asked.
  \item \textbf{Memory}: every \code{new TreeNode} leaks if the harness reuses the
        \code{Solution}; real systems free or pool -- know the ownership model.
\end{itemize}
\end{hardeningbox}

\solvedproblem{c03-dup-subtrees}{Find Duplicate Subtrees}{LeetCode 652}
  {https://leetcode.com/problems/find-duplicate-subtrees/}{LC Medium L5}
  {canonical serialisation, hash-map counting, postorder}

\begin{problemstatement}[Find Duplicate Subtrees]
Return the roots of all duplicate subtrees (two subtrees are duplicate if they have the
same structure \emph{and} same node values). For each repeated shape you may return any
one of its roots. Values in $[-10^{5}, 10^{5}]$, $n \le 5000$.
\end{problemstatement}

\begin{attemptgate}[Attempt it yourself first -- \lv{5}, give yourself 30 minutes]
A fair attempt: find a way to say ``these two subtrees are the same'' cheaply.
\peek{what single string could \emph{fingerprint} a subtree's shape and values?}
\nudge{a null-marked serialisation of each subtree is a canonical name; count names in a
map; emit the root the second time a name repeats.}
\fullhint{Postorder: build \code{key = val + "," + leftKey + "," + rightKey} (with \code{"\#"}
for null), \code{++freq[key]}, and push the root when the count hits exactly $2$. Full build
below.}
\end{attemptgate}

\solstep{Understanding the Problem}
Subtree equality is \emph{structure} $+$ \emph{values}, which is exactly what a null-marked
serialisation encodes (\defref{c03-serialize}, \corref{c03-serialize-bijection}): two
subtrees are identical iff their marked serialisations are identical strings. So the task
becomes ``count equal strings and report the ones seen twice or more.''

\solstep{Key Observations}
\begin{observation}
A postorder assembly -- children's keys first, then the node's -- lets each node's key be
built once from its children in $\bigO{1}$ concatenations.
\end{observation}
\begin{observation}
Null markers are mandatory: without them, a left-only and a right-only single-child node
share a key (\figref{c03-ambig} again), merging distinct shapes.
\end{observation}
\begin{observation}
Pushing the root only when \code{++freq[key] == 2} reports each repeated shape exactly once.
\end{observation}

\solstep{Think It Through With Me}
\begin{thinkbox}
\deadend{Compare every pair of subtrees structurally.}{There are $\bigO{n}$ subtrees and
pairwise comparison is $\bigO{n^2}$ pairs, each $\bigO{n}$ to check -- hopeless.}
\deadend{Fingerprint by the multiset of node values only.}{Two different \emph{shapes} can
share a value multiset (e.g.\ a chain vs a fork on $\{1,2,3\}$); values without structure
over-merge. The key must encode \emph{shape}, which null-marked serialisation does.}
\aha{Use the null-marked preorder key computed bottom-up, hash it, and group by the hash:
equality of keys $\Leftrightarrow$ equality of subtrees (\corref{c03-serialize-bijection}).}
\end{thinkbox}

\solstep{The Approach Ladder}
\begin{center}\renewcommand{\arraystretch}{1.3}
\begin{tabular}{@{}c p{0.225\linewidth} p{0.255\linewidth} p{0.235\linewidth} c@{}}
\toprule
\textbf{\#} & \textbf{Approach} & \textbf{Idea} & \textbf{Time} & \textbf{Verdict} \\
\midrule
1 & \lv{3} pairwise check & compare all subtree pairs & $O(n^3)$ & \tle \\
2 & \lv{4} value-multiset key & hash values only & $O(n^2)$ & \textbf{WA} \\
3 & \lv{5} marked-serial key & shape+value string & $O(n^2)$, $O(n)$ intern'd & \ac \\
\bottomrule
\end{tabular}\end{center}

\solstep{Solution}
Postorder DFS returning the subtree's serialisation string. Maintain
\code{unordered\_map<string,int> freq}. At each node form \code{key}; if
\code{++freq[key]==2} add the node to the answer.

\solstep{Proof of Correctness}
\begin{proofbox}
By \corref{c03-serialize-bijection} the marked-preorder string is a bijection between
subtrees and well-formed strings, so two nodes' subtrees are duplicates iff their keys are
equal. The postorder computes each node's key from its children's keys, which equal their
subtrees by the same fact; hence \code{freq[key]} counts subtrees isomorphic (in value and
structure) to $\subtree(v)$. Emitting $v$ exactly when the count reaches $2$ yields one
representative per equivalence class of size $\ge 2$ and none of size $1$ -- precisely the
required output.
\end{proofbox}

\solstep{Implementation}
\begin{lstlisting}[caption={Find Duplicate Subtrees (LeetCode 652)}, label={lst:c03-dup-subtrees}]
#include <bits/stdc++.h>
using namespace std;

struct TreeNode {
    int val;
    TreeNode *left, *right;
    TreeNode(int x = 0) : val(x), left(nullptr), right(nullptr) {}
};

class Solution {
    unordered_map<string,int> freq;
    vector<TreeNode*> res;
    string key(TreeNode* t) {
        if (!t) return "#";                              // null marker -> shape-sensitive
        string s = to_string(t->val) + "," + key(t->left) + "," + key(t->right);
        if (++freq[s] == 2) res.push_back(t);            // emit each dup once
        return s;
    }
public:
    vector<TreeNode*> findDuplicateSubtrees(TreeNode* root) {
        key(root);
        return res;
    }
};
\end{lstlisting}

\solstep{Trace}
\begin{tracebox}[On $R$ (leaves shared, internal not)]
Leaves: \code{key(4)="4,\#,\#"}, \code{"5,\#,\#"}, \code{"6,\#,\#"} -- all first-time. Node
$2$: \code{"2,4,\#,\#,5,\#,\#"} first time; node $3$: \code{"3,\#,6,\#,\#"} first time;
root $1$: unique. No key reaches count $2$ on $R$ (its three leaf keys are \emph{distinct}
because values differ), so result is empty -- correct, $R$ has no repeated subtree. If two
leaves shared a value, say two \code{7}s, their \code{"7,\#,\#"} key would hit $2$ once and
be reported.
\end{tracebox}

\solstep{Complexity}
\begin{complexitybox}
Naive string keys: node at depth $d$ builds an $\bigO{\text{subtree size}}$ string, total
$\bigO{n^2}$ time and $\bigO{n^2}$ hash storage worst case (a comb tree). Interning each
unique key to a small integer id (see hardening) makes the \emph{distinct} subtree count
the work bound, typically $\bigO{n}$ average. Postorder DFS stack $\bigO{h}$.
\end{complexitybox}

\solstep{Verification and Edge Cases}
\begin{center}\renewcommand{\arraystretch}{1.2}
\begin{tabular}{@{}p{0.30\linewidth} p{0.62\linewidth}@{}}
\toprule
\textbf{Case} & \textbf{What the code does} \\
\midrule
Identical single-node subtrees & \code{"7,\#,\#"} repeats; reported once \\
Whole tree repeated (root counted twice) & only the first node to reach \code{freq==2}
  is added, never the ancestor of an already-duplicate -- LC forbids nested duplicates in
  the answer, which the \code{==2} rule matches \\
All-distinct values, unique shapes & empty result (as on $R$) \\
\bottomrule
\end{tabular}\end{center}

\solstep{Pitfalls and Variants}
\begin{warningbox}
Dropping the null markers (\code{"\#"}) or the separators merges distinct shapes (Traps 2
and 3 of \secref{c03-rebuild-pitfalls} in one). Emit on \code{==2} exactly, not
\code{>=2} (which would list the same shape once per extra copy).
\end{warningbox}

\begin{insightbox}[Key takeaways -- three things to remember]
\begin{enumerate}[nosep, leftmargin=1.6em]
  \item A \emph{canonical serialisation} is a fingerprint: reduce ``are these subtrees
        equal?'' to ``are these strings equal?'' and let a hash map do the grouping.
  \item Shape-sensitivity requires explicit nulls -- the same lesson the (pre,post)
        ambiguity taught us.
  \item ``Report on the transition to count $2$'' is the generic one-representative-per-class
        idiom; it recurs in de-duplication problems across the book.
\end{enumerate}
\end{insightbox}

\begin{hardeningbox}
\begin{itemize}[nosep, leftmargin=1.4em]
  \item \textbf{Quadratic strings}: interning --- map each distinct key string to an
        \code{int} id and represent a subtree by the triple \code{(val, leftId, rightId)}
        in a \code{map<tuple,int>} --- makes each node $\bigO{1}$ expected work, total
        $\bigO{n}$; this is the honest answer to the ``can you do better?'' follow-up.
  \item \textbf{Adversarial depth}: a chain of $5000$ nodes is $\bigO{n}$ recursion; if the
        judge stack is tight, convert the postorder to an explicit stack (\secref{c03-iterative}).
  \item \textbf{Collisions}: if you hash to a fixed-width integer instead of the full
        string, a collision silently merges classes -- prefer exact strings, or the interned
        id, not a lossy hash, when correctness is graded.
\end{itemize}
\end{hardeningbox}

\subsection{Practice Problems}

\begin{stuckbox}
\begin{itemize}[nosep, leftmargin=1.4em]
  \item In reconstruction: did you exclude the root token from both child postorder
        windows? (Off-by-one, \probref{c03-build-in-post}.)
  \item Do you have a value $\to$ inorder-index map, or are you searching each level?
        (Trap 1.)
  \item Does your serialisation write a marker for \emph{every} null and a separator between
        values? (Traps 2--3.)
  \item For subtree hashing: would removing a null marker merge two shapes that differ only
        in a child's side?
\end{itemize}
\end{stuckbox}

\begin{practiceset}
  \practice{Construct Binary Tree from Preorder and Postorder Traversal}{LeetCode 889}
    {https://leetcode.com/problems/construct-binary-tree-from-preorder-and-postorder-traversal/}{LC Medium L3}
    {The pair is \emph{not} unique (\figref{c03-ambig}); LC asks for \emph{any} valid tree,
    so when a node has one child, assign it consistently to the left and advance. Root from
    preorder's front, split by the next preorder token's postorder index.}
  \practice{Serialize and Deserialize BST}{LeetCode 449}
    {https://leetcode.com/problems/serialize-and-deserialize-bst/}{LC Medium L4}
    {A BST's preorder alone deserialises uniquely (no null markers needed) -- bound each
    parse by the current $[\mathit{lo},\mathit{hi}]$ value window. Contrast with the general
    tree needing markers; the Summit LC 297 is exactly this without the BST promise.}
\end{practiceset}

\begin{recallbox}[10-second refresher -- 3.4 Rebuilding and serialising]
Reconstruction needs a \emph{splitter}: (pre,in) or (post,in) give a unique tree because
inorder's root index carves left/right; (pre,post) is ambiguous at one-child nodes. To
round-trip a tree to a string, mark every null and separate every token -- a marked
preorder is a bijection with the tree, and doubles as a subtree fingerprint for hashing.
\selftest{
  \item State \thmref{c03-pair-uniqueness}: which pairs are unique and which is not, and
        give the one-child counterexample.
  \item Write the reconstruction recurrence (\coderef{c03-build}) from memory and mark the
        two window-arithmetic lines.
  \item Why is the value $\to$ inorder-index map needed for $\bigO{n}$ time?
  \item Give the two conditions a serialisation string must satisfy to deserialise to
        exactly one tree.
  \item On LC 652, why does emitting on \code{freq==2} (not \code{>=2}) give one node per
        duplicate class?
}
\end{recallbox}
